\documentclass[aspectratio=169,11pt]{beamer}
\usepackage[utf8]{inputenc}
\usepackage[T1]{fontenc}
\usepackage[brazilian]{babel}
\usepackage{amsmath,booktabs,graphicx,array}
% Figuras pré-compiladas em figuras/. Dados e falas documentados no roteiro.
% Recompilar figuras: powershell -File figuras/compilar-figuras.ps1
% Estética: Singapore, como semanas 1--3. Conteúdo conciso e visual: Mini-ARC.
\usetheme{Singapore}
\title[PHVD-APS -- Apresentação final]{Planejamento de Visitas Domiciliares na APS}
\subtitle{Slides de apoio}
\author{Fabio Cieslak, Tobias Marion e Vitor Feijo}
\date{INF99003 -- Projeto 2 -- Grupo A\\Apresentação final -- 09/10/2026}
\definecolor{ink}{HTML}{202124}
\definecolor{muted}{HTML}{666666}
\definecolor{blue}{RGB}{30,80,160}
\definecolor{green}{RGB}{30,130,60}
\definecolor{red}{RGB}{190,60,60}
\definecolor{light}{HTML}{E8E8E8}
\definecolor{amber}{HTML}{B68432}
\setbeamercolor{normal text}{fg=ink,bg=white}
\setbeamersize{text margin left=8mm,text margin right=8mm}
\setbeamertemplate{navigation symbols}{}
\setbeamercovered{transparent}
\setbeamertemplate{footline}[frame number]
\newcommand{\bigtext}[1]{{\fontsize{24.88}{30}\selectfont #1\par}}
\newcommand{\smallnote}[1]{\par\vspace{2mm}{\fontsize{8}{10}\selectfont\color{muted}#1\par}}
\newcommand{\labeltext}[1]{{\color{muted}\small #1}}
\begin{document}
\section{Apoio}
% Apoio 1
\begin{frame}{Apoio · como medimos atendimento}
\vspace{3mm}
\begin{tabular}{@{}p{.31\textwidth}p{.62\textwidth}@{}}
\textbf{Prioridade a tempo} & Pesos atendidos a tempo / pesos da demanda inicial\\[6mm]
\textbf{Cobertura efetiva} & Visitas concluídas / visitas da demanda inicial\\[6mm]
\textbf{Atraso controlável} & Dias adicionais de atraso, ponderados pelos pesos
\end{tabular}
\vspace{7mm}\par{\large Pesos 1, 3 e 5; só o peso 5 foi atendido a tempo:}
\par\vspace{3mm}{\large $5/9 = 55{,}56\%$ da prioridade}
\smallnote{Prioridade e cobertura multiplicadas por 100. Já vencidas contam como prontas se atendidas no primeiro dia. Pendências permanecem no denominador.}
\end{frame}
% Apoio 2
\begin{frame}{Apoio · 250 pacientes: capacidade e cálculo}
{\large Restinga · 3 sementes · 1 equipe de 300 min/dia\par}
\vspace{5mm}
\begin{columns}[T]
\column{.54\textwidth}\textbf{Visitas no plano inicial}\par\vspace{3mm}
{\small\begin{tabular}{@{}lrr@{}}
\toprule Método & 4 meses & 12 meses\\\midrule
Principal & 249--250 & 250\\
Vizinho & 241--243 & 250\\
Urgência & 246--248 & 250\\\bottomrule
\end{tabular}}
\column{.43\textwidth}\textbf{Cálculo da principal}\par\vspace{3mm}
{\small\begin{tabular}{@{}lr@{}}
\toprule Ausência & Média acumulada\\\midrule
0\% & 22,14 s\\
5\% & 116,20 s\\
10\% & 151,78 s\\\bottomrule
\end{tabular}}
\end{columns}
\vspace{5mm}\labeltext{Com 240 min/dia, algumas visitas isoladas excediam a jornada.}
\smallnote{Capacidade: mínimo--máximo entre sementes, de 250 visitas. Cálculo: médias das execuções longas por taxa. Limite comum de 261 dias úteis; feriados não retirados.}
\end{frame}
% Apoio 3
\begin{frame}{Apoio · protocolo e função de busca}
\begin{columns}[T]
\column{.47\textwidth}\textbf{Fatorial}\par\vspace{3mm}
\begin{tabular}{@{}ll@{}}
Pacientes & 15 / 30 / 45 / 90\\[2mm]
Área & 0,5 / 1 / 2$\times$\\[2mm]
Dias úteis & 5 / 10 / 22\\[2mm]
Já vencidas & 0 / 25 / 50\%\\[2mm]
Falha & 0 / 5 / 10\%
\end{tabular}
\column{.5\textwidth}\textbf{Busca entre rotas}\par\vspace{3mm}
{\small $J=T_{\mathrm{viagem}}+12D_w-24P_{\mathrm{pronta}}$}\par
\vspace{4mm}\labeltext{$D_w$: atraso controlável ponderado\\$P_{\mathrm{pronta}}$: pontos de prioridade pronta}
\par\vspace{4mm}\labeltext{Primeira busca: peso da resposta pronta = 0\\Segunda busca: peso = 24}
\end{columns}
\smallnote{324 cenários; 2.916 simulações; 972 pares por referência. OSRM a pé. Pesos em minutos equivalentes; até 80 movimentos por busca.}
\end{frame}
% Apoio 4
\begin{frame}{Apoio · reprodução dos três estudos}
\relax
{\large OSRM local com OpenStreetMap e perfil de caminhada\par}
\vspace{6mm}
\begin{tabular}{@{}ll@{}}
Porto Alegre & \texttt{npm run citywide:report}\\[4mm]
Fatorial & \texttt{npm run factorial}\\[2mm]
 & \texttt{npm run factorial:report}\\[4mm]
250 pacientes & \texttt{npm run long-horizon}
\end{tabular}
\vspace{6mm}\par\labeltext{Comandos com \texttt{\mbox{-}\mbox{-}prefix experimentos} e \texttt{OSRM\_BASE\_URL} definido}
\smallnote{Compilar core, server e experimentos antes. Sequência completa, capacidade e verificação das matrizes no roteiro e em experimentos/README.md.}
\end{frame}
\end{document}
